C: Das wird ein Spaß.
R: ich mache mal einen commit und pushe den.

Hab jetzt die Datei geändert und gespeichert. Dann "git add *" um alles zu adden oder "git add ." um das aktuelle Verzeichnis zu adden oder halt "git add main.tex" um explizit zu adden.
Dann "git commit -m 'mein erster commit'" und dann "git push". Mal gucken obs klappt :)

C: Ich habe homebrew und gh installiert.
C: ´git add´ ist das gleiche wie stagen mit dem Plus links bei VSCode.

C: gh issue close 2
C: gh pr merge 1

C: `git pull --rebase` ist, wie der Name schon sagt, auch eine Möglichkeit, um die Änderungen vom Remote-Repository zu holen und gleichzeitig die eigenen Änderungen darauf anzuwenden. Das kann helfen, Konflikte zu vermeiden. 

WICHTIG: bei `gh pr merge 1` muss man DIREKT danach `git pull` machen, sonst hat man die Änderungen nicht lokal und kann nicht weiterarbeiten.

R: um im lokalen repo zu checkn ob es remote änderungen gab muss zuerst `git fetch` ausgeführt werden. dann kann mit `git status` nachgeschaut werden, ob es commits gibt, die gepullt werden können/müssen.

C: Autofetch: über Cmd + Shift + P -> User Preferences und dann über git.autofetch auf true setzen. Dann wird automatisch alle 3 Minuten ein fetch gemacht. Man kann auch die Period ändern.

C: Wenn beide zeitgleich an einer Datei gearbeitet haben, und zB Renas gepusht hat, dann muss ich meine Änderungen stashen, dann pullen, dann über den Merge-Editor in VSCode die Änderungen zusammenführen und dann committen. Danach kann ich pushen (nach add und commit).